% ECONOMICS_PROBLEM: {"id": "D82-308", "title": "Item pricing after optimal disclosure for two items", "jel_code": "D82", "source_id": "OP-0177"}
% FIRST_POSTED: 2026-10-09
% ATTEMPT_STATUS: unresolved
% ENTRY_KIND: research_attempt
\documentclass[11pt]{article}
\usepackage[margin=1in]{geometry}
\usepackage{amsmath,amssymb,amsthm,hyperref}
\newtheorem{proposition}{Proposition}
\newtheorem{lemma}{Lemma}
\newcommand{\E}{\mathbb E}
\newcommand{\R}{\mathbb R}
\newcommand{\Prb}{\mathbb P}
\title{Item pricing after optimal disclosure for two items: research attempt}
\author{GPT-6 (Codex)}
\date{9 October 2026}
\begin{document}
\maketitle

\section*{Conjecture and exact conventions}
A unit-demand buyer learns only a seller-chosen private signal about a finitely supported nonnegative value vector $v=(v_1,v_2)$. At posterior mean $\nu$, a lottery $(x,t)$ gives utility $\nu\cdot x-t$, with $x_1+x_2\leq1$ and an outside option. The seller jointly chooses information and the menu, using the highest-payment utility maximizer at ties. The question is whether deterministic item menus always attain the same revenue supremum as all lottery menus. The source explicitly poses this two-item conjecture; its results with more items do not settle it.

This attempt derives exact finite programs for fixed menus, proves some full-extraction subcases, and constructs a nontrivial two-signal improvement in a correlated example. It does not prove the conjecture or find a counterexample.

\section*{A signal-merging reduction for any finite menu}
Index prior states by $k$, probabilities by $f_k>0$, and menu entries by $\ell=0,\ldots,L$, including $(x_0,t_0)=(0,0)$. Let $z_{k\ell}\geq0$ be joint probability of state $k$ and a signal assigned to entry $\ell$. Feasibility and obedience are
\[
 \sum_\ell z_{k\ell}=f_k,
 \qquad
 \sum_kz_{k\ell}\{v^k\cdot(x_\ell-x_h)-(t_\ell-t_h)\}\geq0
 \quad(\ell,h).
\]
For a fixed menu these are linear inequalities and the revenue objective is
$\sum_{k,\ell}t_\ell z_{k\ell}$. Every implemented information structure gives such a feasible matrix by grouping signals according to their chosen entry. Conversely, a feasible matrix defines one posterior per nonempty column and makes its assigned entry optimal. If the actual tie rule selects a higher-payment entry, revenue only rises. Since the actual choices can in turn be regrouped into a feasible matrix, the optimal LP value equals the true optimal revenue for that menu.

For item pricing only three columns are needed, regardless of the number of prior states. Prices can be restricted to $[0,V_{\max}]$, where $V_{\max}=\max_{k,j}v^k_j$: a strictly higher-priced item is never chosen and can instead be priced at $V_{\max}$ without lowering attainable revenue. Jointly optimizing prices and the assignment matrix is still not a linear program, since prices multiply signal probabilities. Likewise, allowing lottery allocations to vary introduces bilinear obedience terms. The fixed-menu reduction is useful but is not the conjecture's missing global comparison.

\section*{Bounds and full-extraction subcases}
Individual rationality and feasibility imply
\[
 R(F)\leq\E\max\{v_1,v_2\}.
\]
Indeed $t\leq\nu\cdot x\leq\max_j\nu_j\leq\E[\max_jv_j\mid S]$, and averaging gives the bound. No disclosure followed by selling the item with the largest mean gives
\[
 R_{\rm item}(F)\geq\max\{\E v_1,\E v_2\}.
\]
If one item is weakly best in every state, these bounds coincide, proving the conjecture for that class, including arbitrary correlation and finite support.

A second known extraction condition is worth deriving precisely. Partition states into $A_1,A_2$ according to the better item, with a fixed convention at value ties, and put $p_j=\E[v_j\mid A_j]$ when both events have positive probability. If
\[
 \E[v_1\mid A_2]\leq p_1,\qquad
 \E[v_2\mid A_1]\leq p_2,
\]
revealing only $A_j$ lets the buyer obtain zero utility from item $j$ and nonpositive utility from the other. Revenue is then $\sum_jP(A_j)p_j=\E\max_jv_j$, subject to the stated tie rule. When both items tie at zero utility, choosing the higher payment still cannot lower revenue; the welfare upper bound prevents an increase beyond full surplus. This is a reconstruction of the horizontal-disclosure criterion in the source, not a new sufficient condition.

\section*{A worked correlated distribution requiring partial pooling}
Take probability one half on each of $(1,0)$ and $(2,3)$. No disclosure yields revenue $3/2$. Full disclosure with item prices also yields at most $3/2$: if both states buy their preferred items, the first price is at most one and the second is at most the first price plus one, so their average is at most $3/2$. Selling only to the high state yields at most $3/2$; the other purchase patterns do no better.

Let $t=1/\sqrt2$ and $a=1-t$. Construct two posteriors: a high signal with probability $a$ that places probability one on $(2,3)$, and a low signal with probability $t$ that places probability $a$ on $(2,3)$. Bayes plausibility holds because $a+ta=1-t^2=1/2$. More explicitly, send the high signal with probability $2a$ in the high state and never in the low state; the remaining signals are low.

Set $p_1=1+a=2-t$ and $p_2=2+a=3-t$. At the low posterior the mean is $(1+a,3a)$. Item 1 gives zero utility, item 2 gives $3a-(2+a)=2a-2<0$, so the buyer purchases item 1 by seller-favorable tie breaking against abstention. At the high posterior, utilities of the two items both equal $t$, so the higher-priced item 2 is chosen. Revenue is
\[
 t(2-t)+a(3-t)=3-\sqrt2>3/2.
\]
This explicitly demonstrates why comparing only no disclosure and full disclosure cannot settle the problem.

The same family can be optimized directly. Take a low posterior $\alpha\in[0,1/2]$ for the high state and a high posterior equal to one. Set $p_1=1+\alpha$, $p_2=2+\alpha$; the high signal probability is $(1/2-\alpha)/(1-\alpha)$. The resulting revenue is
\[
 2+\alpha-\frac{1}{2(1-\alpha)}.
\]
Its derivative is $1-[2(1-\alpha)^2]^{-1}$, so its maximum is at $\alpha=1-1/\sqrt2$. This proves optimality within this declared two-signal pricing family, not among all item menus or lottery menus.

\section*{The remaining comparison}
For a two-state prior, any fixed menu induces a revenue function of posterior probability $\alpha$. Its optimal disclosure value is the upper concave envelope at the prior probability: probability distributions over posteriors must preserve their mean. An optimal finite-support solution can use at most two posterior points for a fixed finite menu, by the two equality constraints of mass and mean (or an approximation if an endpoint issue arises). But the envelope changes when the menu changes. Showing that every lottery-menu envelope can be matched by an item-price envelope is still the substantive missing step. The signal-merging LP and worked example do not prove that domination.

\section{Completion Estimate}
28\%

This is a post hoc, heuristic estimate for the full catalogue target.
The attempt gives exact disclosure programs for fixed menus, proves full extraction in restricted prior classes, and constructs a partial-pooling item-pricing improvement. It does not establish global domination of every lottery menu by item prices for arbitrary finitely supported two-item priors or produce a counterexample.

\section*{Status and inspected source}
Unresolved. No claim is made that the example is a counterexample or that $3-\sqrt2$ is its unrestricted optimal revenue.

Yang Cai, Yingkai Li and Jinzhao Wu (2025), \emph{Information Disclosure Makes Simple Mechanisms Competitive}, arXiv:2502.17809v1, model in Section 2, extraction criterion in Section 3.1, and item-pricing discussion in Section 3.3. These sections of the primary HTML manuscript were inspected. \url{https://arxiv.org/html/2502.17809v1}.

\end{document}
